\documentclass[11pt]{article}

\usepackage[margin=1in]{geometry}
\usepackage{amsmath}
\usepackage{amssymb}
\usepackage{booktabs}
\usepackage{hyperref}
\usepackage{enumitem}

\title{DynamicsCPT: Dynamics So Far}
\author{}
\date{\today}

\begin{document}
\maketitle

\section{Purpose}
This note summarizes the dynamics model currently implemented in the repository. It is not a final theory document. Instead, it captures the present working formulation, the main numerical choices, and the diagnostics that are being used to evaluate whether the collision logic behaves consistently.

\section{State Variables}
Each particle $i$ carries
\begin{align*}
    \mathbf{x}_i &= (x_i, y_i), \\
    \mathbf{p}_i &= (p_{x,i}, p_{y,i}), \\
    m_i &> 0, \\
    r_i &> 0.
\end{align*}
The velocity is recovered from momentum by
\[
    \mathbf{v}_i = \frac{\mathbf{p}_i}{m_i}.
\]

The simulation advances with an outer time step $\Delta t$ and a substep count $N_s$, so the nominal internal step is
\[
    \delta t = \frac{\Delta t}{N_s}.
\]
At the time of writing, the base configuration uses $\Delta t = 0.005$ and $N_s = 5$, so $\delta t = 0.001$ in the default setup.

\section{Geometric Contact Model}
For two disks with centers $\mathbf{x}_i$, $\mathbf{x}_j$ and radii $r_i$, $r_j$, let
\[
    \mathbf{d}_{ij} = \mathbf{x}_j - \mathbf{x}_i,
    \qquad
    d_{ij} = \|\mathbf{d}_{ij}\|.
\]
Contact occurs when
\[
    d_{ij} < r_i + r_j.
\]
When there is contact, the code computes a contact normal
\[
    \mathbf{n}_{ij} =
    \begin{cases}
        \dfrac{\mathbf{d}_{ij}}{d_{ij}}, & d_{ij} > 0, \\
        (1,0), & d_{ij} = 0,
    \end{cases}
\]
and an overlap area $A_{ij}$ from the exact two-circle intersection formula. The same general approach is used for wall contact, with a wall normal replacing the pair normal.

The overlap area is the main geometric quantity that drives the collision response. In other words, the current model is not event-based in the hard-sphere sense; it is overlap-driven.

\section{Response Laws}
\subsection{Force-form and acceleration-form coefficients}
The code supports two equivalent ways to specify response strength:
\begin{itemize}[leftmargin=1.5em]
    \item a force-per-area coefficient $k$, or
    \item an acceleration-per-area coefficient $a_A$.
\end{itemize}

For wall contact, the effective force magnitude is
\[
    F = k A
\]
in force form, or
\[
    F = m_i a_A A
\]
in acceleration form.

For particle-particle contact, the acceleration-form implementation uses reduced mass
\[
    \mu_{ij} = \frac{m_i m_j}{m_i + m_j},
\]
so the force magnitude becomes
\[
    F = \mu_{ij} a_A A.
\]

\subsection{Impulse over one substep}
The response is applied over the substep as an impulse
\[
    J = F \, \delta t.
\]
For a pair contact, the momentum update has the form
\begin{align*}
    \mathbf{p}_i^{\,new} &= \mathbf{p}_i - J \mathbf{n}_{ij}, \\
    \mathbf{p}_j^{\,new} &= \mathbf{p}_j + J \mathbf{n}_{ij}.
\end{align*}
For wall contact, the particle receives the wall-normal impulse directly.

\section{Two Active Dynamics Variants}
\subsection{Base model}
The base class contains the common geometry, bookkeeping, momentum and energy diagnostics, and collision trace export. It also supports optional wall contact and optional bonded interactions.

\subsection{FieldBase}
The currently active scenario builder constructs a \texttt{FieldBase} object. In this formulation, particles are treated as field sources and the overlap area acts as a coupling measure. Important features are:
\begin{itemize}[leftmargin=1.5em]
    \item overlap-driven impulses rather than event collisions,
    \item per-particle or per-pair response coefficients,
    \item inbound and outbound contact accounting,
    \item turning-area diagnostics for off-center collisions,
    \item internal scalar bookkeeping to track changes in the sum of speed magnitudes.
\end{itemize}

The code distinguishes inbound contact ($v_{rel,n} < 0$) from outbound contact ($v_{rel,n} \ge 0$), where
\[
    v_{rel,n} = (\mathbf{v}_j - \mathbf{v}_i)\cdot \mathbf{n}_{ij}.
\]
This inbound/outbound split is used in the diagnostics and in the outbound release logic.

\subsection{AccelBase}
\texttt{AccelBase} is a related variant that also uses overlap-driven dynamics, but its bookkeeping is organized slightly differently. In particular, it keeps a \texttt{bookkeeping\_internal} term and allows free-flight settling steps after contact to push the scalar accounting toward closure.

\section{Position Update}
Once the momentum updates for a substep are complete, the particle positions are advanced explicitly:
\[
    \mathbf{x}_i^{\,new} = \mathbf{x}_i + \mathbf{v}_i^{\,new}\,\delta t.
\]
This is therefore an explicit time-marching scheme with contact impulses applied inside each substep.

\section{Adaptive Refinement Near Contact}
The base integrator can refine a substep when contact is active and the current speed would traverse too much contact span in one step. In effect, the code estimates a required number of refined steps and recursively subdivides the current substep if needed.

The current logic uses a minimum-collision-steps target to bound how much relative motion can occur over one contact span. This is meant to reduce under-resolution during fast impacts and to make the area-driven impulse more stable.

\section{Bonded Interactions}
Optional bonded interactions are also present in the code. A bonded pair can apply
\begin{itemize}[leftmargin=1.5em]
    \item a spring-like contribution proportional to extension, and
    \item a damping-like contribution proportional to normal relative velocity.
\end{itemize}
These act over the same substep interval and are added before the contact position writeback. Bonding is not the main focus of the current two-particle overlap study, but it is already part of the framework.

\section{Diagnostics and Invariants}
Several diagnostics are exported at the end of a run.

\subsection{Vector momentum}
The total vector momentum is
\[
    \mathbf{P} = \sum_i \mathbf{p}_i.
\]
The code tracks the difference between initial and final vector momentum as a direct conservation check.

\subsection{Kinetic energy}
The kinetic energy is
\[
    K = \sum_i \frac{\|\mathbf{p}_i\|^2}{2m_i}.
\]
This is not guaranteed to be exactly conserved in the present overlap-driven formulation, so it is monitored as a diagnostic rather than imposed as a strict invariant.

\subsection{Scalar momentum plus internal term}
An important repository-specific quantity is the scalar total
\[
    S = \sum_i \|\mathbf{p}_i\| + P_{\text{internal}},
\]
or, in the speed-based variants,
\[
    S = \sum_i \|\mathbf{v}_i\| + P_{\text{internal}}.
\]
The exact meaning of the internal term is model-dependent, but operationally it is a bookkeeping quantity that compensates for scalar-magnitude changes introduced by the overlap response. The code reports the final scalar error as
\[
    S_{\text{final}} - S_{\text{initial}}.
\]
This scalar closure metric is currently one of the main health checks for the collision model.

\subsection{Overlap and impulse symmetry}
The runs also record:
\begin{itemize}[leftmargin=1.5em]
    \item maximum inbound and outbound overlap area,
    \item maximum inbound and outbound impulse,
    \item cumulative area, area-time, and impulse totals,
    \item turning-area measures for glancing contacts.
\end{itemize}
These diagnostics are intended to show whether the inbound and outbound stages are balanced, and whether off-axis turning correlates with scalar error.

\section{Current Two-Particle OCL Study}
The main batch study in \texttt{run\_two\_particle\_ocl.py} varies the initial $y$-position of particle 1 while leaving particle 0 on the horizontal axis. This creates a family of central-to-glancing collisions.

In the current sweep:
\begin{itemize}[leftmargin=1.5em]
    \item particle 0 starts near $(1.1, 0)$ with leftward velocity,
    \item particle 1 starts near $(-1.1, y)$ with rightward velocity,
    \item the particle 1 vertical offset $y$ is stepped across the study set,
    \item each run exports a full contact trace and a final summary.
\end{itemize}

This study is being used to answer two practical questions:
\begin{enumerate}[leftmargin=1.5em]
    \item How does glancing geometry change overlap, turning, and scalar closure?
    \item How much extra computation is caused by active collision logic relative to free flight?
\end{enumerate}

\section{Measured Processing Time}
The repository now includes frame-timing instrumentation in the visual driver. For each outer simulation frame, the code records:
\begin{itemize}[leftmargin=1.5em]
    \item simulation wall time per frame,
    \item total frame wall time,
    \item render wall time,
    \item active contact counts.
\end{itemize}
This allows direct comparison between free-flight frames and collision-active frames. The goal is to identify how much the contact logic slows the run when overlaps, turning, or refinement become significant.

\section{Interpretation of the Present Model}
So far, the project should be understood as an overlap-driven contact dynamics experiment rather than a finished collision theory. The current implementation is trying to achieve three things at once:
\begin{itemize}[leftmargin=1.5em]
    \item geometrically meaningful contact response through overlap area,
    \item numerically stable time marching through substeps and optional refinement,
    \item near-closure of vector momentum, scalar momentum bookkeeping, and useful symmetry diagnostics.
\end{itemize}

This makes the code especially useful for comparative studies, because each design change can be tested against the same exported metrics.

\section{Open Questions}
The most important unresolved questions appear to be:
\begin{itemize}[leftmargin=1.5em]
    \item whether the scalar internal bookkeeping has a clean physical interpretation or is mainly a numerical closure device,
    \item how outbound release should scale for turning collisions,
    \item how sensitive the results are to $\delta t$, substeps, and contact-refinement settings,
    \item whether the field-based formulation converges toward a stable continuum or hard-contact limit,
    \item how strongly collision geometry correlates with both scalar error and runtime cost.
\end{itemize}

\section{Summary}
At present, the repository implements a two-dimensional particle dynamics framework with overlap-area contact forces, optional adaptive refinement, optional bonding, and detailed trace export. The main working hypothesis is that overlap geometry can drive a useful collision response while scalar bookkeeping tracks the non-vector part of the exchange. The next phase is to use the exported study metrics and frame-timing data to determine which parts of the collision logic are physically helpful, which are only numerically compensating, and which are the main source of runtime slowdown.

\end{document}
